\documentclass[10pt]{article}
\usepackage[T1]{fontenc}
\usepackage[utf8]{inputenc}
\usepackage[french]{babel}
\usepackage{amsmath,amssymb,mathrsfs,mathpazo}
\usepackage{geometry}
\geometry{a4paper,margin=2cm}
\usepackage{graphicx,array,tabularx,tabularray,colortbl,multirow,multicol,booktabs,enumitem,pifont,ulem,fancybox,diagbox,cancel}
\usepackage{pstricks,pstricks-add,pst-plot,pst-tree,pst-node,pst-text,pst-func,pst-math,pst-bspline,pst-3dplot,pst-eucl,pst-solides3d}
\let\pscalculate\undefined
\usepackage{pst-bezier}
\usepackage{tikz,pgfplots,tkz-tab}
\usetikzlibrary{arrows,arrows.meta,patterns,calc,positioning,trees,decorations.markings,shapes.geometric,shapes.misc,shapes.arrows,intersections,decorations.pathmorphing,decorations.pathreplacing,backgrounds,shadows,babel}
\pgfplotsset{compat=1.18}
\usepgfplotslibrary{fillbetween}
\usepackage[np]{numprint}
\usepackage{listings,makecell,pas-tableur,twoopt}
\let\barre\undefined
\usepackage[pstricks]{tablvar}
\definecolor{darkgreen}{RGB}{0,100,0}
\providecommand{\red}{\color{red}}
\providecommand{\blue}{\color{blue}}
\providecommand{\green}{\color{green}}
\providecommand{\black}{\color{black}}
\providecommand{\white}{\color{white}}
\providecommand{\gray}{\color{gray}}
\providecommand{\cyan}{\color{cyan}}
\providecommand{\magenta}{\color{magenta}}
\providecommand{\yellow}{\color{yellow}}
\providecommand{\violet}{\color{violet}}
\providecommand{\orange}{\color{orange}}
\providecommand{\euro}{\texteuro}
\providecommand{\EUR}{\texteuro}
\providecommand{\ieme}{\textsuperscript{e}}
\providecommand{\degres}{\ensuremath{^\circ}}
\providecommand{\degree}{\ensuremath{^\circ}}
\providecommand{\vv}[1]{\overrightarrow{#1}}
\providecommand{\Coord}[1]{#1}
\providecommand{\decosix}{}
\providecommand{\starredbullet}{\ensuremath{\bullet}}
\pagestyle{empty}
\setlength{\parindent}{0pt}
\usepackage{adjustbox}\DeclareUnicodeCharacter{2236}{:}\DeclareUnicodeCharacter{2190}{\ensuremath{\leftarrow}}
\providecommand{\R}{\mathbb{R}}
\providecommand{\N}{\mathbb{N}}
\providecommand{\Z}{\mathbb{Z}}
\providecommand{\Q}{\mathbb{Q}}
\providecommand{\D}{\mathbb{D}}
\providecommand{\C}{\mathbb{C}}
\providecommand{\e}{{\mathrm{e}}}
\providecommand{\vect}[1]{\overrightarrow{#1}}
\providecommand{\vectt}[1]{\overrightarrow{\mathrm{#1}}}
\providecommand{\vvt}[1]{\overrightarrow{\mathrm{#1}}}
\providecommand{\np}[2][]{{#2~#1}}
\providecommand{\vv}[1]{{\overrightarrow{#1}}}
\renewcommand{\d}{{\mathrm{d}}}
\providecommand{\degres}{{^\circ}}
\providecommand{\degre}{{^\circ}}
\providecommand{\blue}{\color{blue}}
\providecommand{\red}{\color{red}}
\providecommand{\enskip}{\quad}
\providecommand{\llbracket}{\lbrack\!\lbrack}
\providecommand{\rrbracket}{\rbrack\!\rbrack}
\providecommand{\euro}{{\text{€}}}
\providecommand{\no}{{\text{n°}}}
\providecommand{\interval}[2]{{[#1;#2]}}
\providecommand{\intervFO}[2]{{[#1;#2[}}
\providecommand{\dotfill}{\ldots}

\providecommand{\ds}{\displaystyle}
\providecommand{\textding}[1]{#1}
\providecommand{\Oij}{\(\mathrm{O};\vec{\imath},\vec{\jmath}\)}
\providecommand{\Oijk}{\(\mathrm{O};\vec{\imath},\vec{\jmath},\vec{k}\)}
\providecommand{\Ouv}{\(\mathrm{O};\vec{u},\vec{v}\)}
\providecommand{\intervFF}[2]{\ensuremath{[#1;#2]}}
\renewcommand{\euro}{\ifmmode\text{\texteuro}\else\texteuro\fi}
\DeclareUnicodeCharacter{00B0}{\ifmmode^{\circ}\else\textdegree\fi}
\DeclareUnicodeCharacter{20AC}{\ifmmode\text{\texteuro}\else\texteuro\fi}
% Énoncés : Éducation nationale. Transcription et corrigés : APMEP, Denis Vergès et contributeurs. Découpage et classement : CoopMaths, Sébastien Lozano, Éric Elter et contributeurs. Adaptation de la présentation : Maths Explorer.
% CC BY-SA : https://coopmaths.fr/www/about/
% https://coopmaths.fr/alea/static/bac/2023/tex/bac_2023_03_sujet1_asie_2.tex
\begin{document}
\noindent\textit{APMEP / CoopMaths. CC BY-SA. Présentation adaptée par Maths Explorer.}\par\medskip
\medskip

On considère le cube ABCDEFGH qui est représenté ci-dessous.

\psset{unit=1cm,arrowsize=2pt 3}
\begin{pspicture}(9.3,9)
\pspolygon(0.5,1.3)(5.5,0.6)(5.5,6.3)(0.5,7)
\uput[dl](0.5,1.3){A} \uput[d](5.5,0.6){B} \uput[u](5.5,6.3){F} \uput[l](0.5,7){E} 
\psline(5.5,0.6)(8.2,1.7)(8.2,7.4)(5.5,6.3)
\uput[r](8.2,1.7){C} \uput[ur](8.2,7.4){G} 
\psline(8.2,7.4)(3.7,8.1)(0.5,7)
\uput[u](3.7,8.1){H}
\psline[linestyle=dashed](0.5,1.3)(3.7,2.4)(3.7,8.1)
\psline[linestyle=dashed](3.7,2.4)(8.2,1.7)
\uput[d](3.7,2.4){D}
\end{pspicture}


Dans le repère orthonormé $\left(\text{A}~;~\vect{\text{AB}}~;~\vect{\text{AD}}~;~\vect{\text{AE}}\right)$, on considère les points M, N et P de coordonnées :
\begin{center}M$\left(1~;~1~;~\dfrac34\right)$, \qquad N$\left(0~;~\dfrac12~;~1\right)$, \qquad P$\left(1~;~0~;~-\dfrac54\right)$\end{center}

Dans cet exercice, on se propose de calculer le volume du tétraèdre FMNP.

\medskip

\begin{enumerate}
\item Donner les coordonnées des vecteurs $\vect{\text{MN}}$ et $\vect{\text{MP}}$.
\item Placer les points M, N et P sur le cube ci-dessus.
\item Justifier que les points M, N et P ne sont pas alignés. 

Dès lors les trois points définissent le plan (MNP).
\item 
	\begin{enumerate}
		\item Calculer le produit scalaire $\vect{\text{MN}} \cdot \vect{\text{MP}}$, puis en déduire la nature du triangle MNP.
		\item Calculer l'aire du triangle MNP.
	\end{enumerate}
\item 
	\begin{enumerate}
		\item Montrer que le vecteur $\vect{n}(5~;~-8~;~4)$ est un vecteur normal au plan (MNP).
		\item En déduire qu'une équation cartésienne du plan (MNP) est $5x - 8y + 4z = 0$.
	\end{enumerate}	
\item On rappelle que le point F a pour coordonnées F(1~;~0~;~1). 

Déterminer une représentation paramétrique de la droite $d$ orthogonale au plan (MNP) et passant par le point F.
\item On note L le projeté orthogonal du point F sur le plan (MNP). 

Montrer que les coordonnées du point L sont: L$\left(\dfrac47~;~\dfrac{24}{35}~;~\dfrac{23}{35}\right)$.
\item Montrer que FL $= \dfrac{3\sqrt{105}}{35}$ puis calculer le volume du tétraèdre FMNP. 

On rappelle que le volume V d'un tétraèdre est donné par la formule :

\begin{center}$V = \dfrac13 \times  \text{aire d'une base} \times \text{hauteur associée à cette base}$.\end{center}

\end{enumerate}

\medskip

\end{document}
